In the last two chapters, we looked at the support for monitors provided by
SCL, namely locks and conditions.  In~this chapter, we look at similar
mechanisms provided by the Java Virtual Machine (JVM).  The fact that these
mechanisms are built in to the JVM makes them more efficient: the SCL
implementations require another layer of interpretation.  However, they have
disadvantages that make them harder to use: we will discuss these shortly.

As with SCL monitors, there are two aspects to JVM monitors: mutual exclusion,
and a wait/signal mechanism.  We start with mutual exclusion in the next
section, and then consider the wait/signal mechanism.

%%%%%%%%%%%%%%%%%%%%%%%%%%%%%%%%%%%%%%%%%%%%%%%%%%%%%%%

\section{Mutual Exclusion Using JVM Monitors}

The mechanism the JVM provides to ensure mutual exclusion is the
\SCALA{synchronized} method.  
%
If |o| is a reference object and \SCALA{e} is an expression or command, then
|o.synchronized{e}| acts much like \SCALA{e}, except that at most
one thread can be active within the \SCALA{synchronized} expressions of the
object~|o| at any time.  When an object is used like this, we will refer to it
as a \emph{monitor}. 

If a thread tries to execute a \SCALA{synchronized} expression while another
thread is within a |synchronized| block on the same monitor, then the former
thread has to wait until the latter leaves its block.  In effect, the monitor
provides a lock: when a thread enters a |synchronized| block, it has to obtain
the lock; when it leaves the block, it releases the lock.  Thus, this can be
used to ensure that executions of code within |synchronized| blocks do not
interfere with one another.

Note that the monitor must be a \emph{reference} object (a subclass of
|AnyRef|; see Scala box~\ref{sb:reference-types}), not, for example, an |Int|.

The expression |synchronized{e}| is shorthand for |this.synchronized{e}|,
i.e.~where the current object is used for the monitor.

Figure~\ref{fig:counter-monitor} gives a straightforward implementation of a
counter object, implemented using a JVM monitor.  Each operation is wrapped in
a |synchronized| block.  The |Counter| object itself acts at the monitor in
this case.  At most one thread at a time can be inside any of the three
|synchronized| blocks.

%%%%%

\begin{figure}
\begin{scala}
object Counter{
  private var x = 0
  def inc() = synchronized{ x += 1 }
  def dec() = synchronized{ x -= 1 }
  def get = synchronized{ x }
}
\end{scala}
\caption{A counter implemented using a JVM monitor.}
\label{fig:counter-monitor}
\end{figure}

%%%%%


\begin{figure}
\begin{scala}
class MonitorTotalQueue[T] extends TotalQueue[T]{
  /** The current state of the queue. */
  private val queue = new Queue[T]

  /** Enqueue x. */
  def enqueue(x: T) = synchronized{ queue.enqueue(x) }

  /** Try to dequeue.  Return None if the queue is empty. */
  def dequeue(): Option[T] = synchronized{ 
    if(queue.isEmpty) None else Some(queue.dequeue()) 
  }
}
\end{scala}
\caption{A total queue implemented using a JVM monitor.}
\label{fig:totalQueue-JVM-monitor}
\end{figure}

JVM monitors can also be used to implement total datatypes.  For example,
Figure~\ref{fig:totalQueue-JVM-monitor} gives a straightforward implementation
of a total queue. 

\begin{instruction}
Study the implementation.
\end{instruction}

%%%%%

JVM monitors are re-entrant, like SCL |Lock|s.  For example, consider an
object of the following form.
%
\begin{mysamepage}
\begin{scala}
object Obj{
  def f(x: Int) = synchronized{ ...; g(y); ... }
  def g(y: Int) = synchronized{ ... }
}
\end{scala}
\end{mysamepage}
%
If a thread is inside~|f|, it holds the lock on the monitor.  When it
calls~|g|, it obtains the lock for a second time.  When |g| returns, it still
holds the lock once.  Only when |f| returns can another thread obtain the
lock.

Correct use of |synchronized| blocks avoids problems with cache consistency.
When a thread leaves a |synchronized| block, it \emph{publishes} the values in
its cache, making them available to other threads.  When a thread enters a
|synchronized| block, it \emph{subscribes} to all variables published by
other threads on the monitor.  This means that if all accesses to a variable
are protected by |synchronized| blocks on the same monitor, as in
Figure~\ref{fig:counter-monitor} and~\ref{fig:totalQueue-JVM-monitor}, then
any read of the variable is guaranteed to see the most-recent value written by
another operation.  Further, compiler optimisations must respect this
semantics.  This means that if you are using |synchronized| blocks correctly
to protect against data races, you do not need to worry about memory
consistency.

The examples we have seen so far have used coarse-grained locking, with a
single monitor object.  However, it is also possible to use finer-grained
locking, with multiple monitors, so long as operations using different
monitors do not interfere.

For example, Figure~\ref{fig:StripedCounterArray-JVM} uses striped locking to
protect an array of counters, in a very similar style to the implementation in
Figure~\ref{fig:StripedCounterArray} (and extending the |CounterArray|
abstract class from Figure~\ref{fig:CounterArray}). As with the earlier
example, each monitor |locks(j)| protects all counters |a(i)| such that $\sm i
\% \sm{stripes} = \sm{j}$, so each |a(i)| is protected by
\SCALA{locks(i\%stripes)}.  Recall that a monitor can be any reference object,
i.e.~any subclass of |AnyRef|.  It makes sense here to use the simplest such
class, i.e.~|AnyRef| itself.
%%%%%

\begin{figure}
\begin{scala}
class StripedCounterArray(n: Int, stripes: Int) 
    extends tacp.locks.CounterArray(n){
  require(stripes > 0)

  /** Locks.  £locks(j)£ protects all £a(i)£ such that £i\%stripes $=$ j£. */
  private val locks = Array.fill(stripes)(new AnyRef)

  def inc(i: Int) = locks(i%stripes).synchronized{ a(i) += 1 }

  def dec(i: Int) = locks(i%stripes).synchronized{ a(i) -= 1 }

  def get(i: Int) = locks(i%stripes).synchronized{ a(i) }
}
\end{scala}
\caption{An array of counters using striped locking.}
\label{fig:StripedCounterArray-JVM}
\end{figure}
 

%%%%%


\begin{figure}
\begin{scala}
class MonitorShardedSet[A](shards: Int) extends Sharding[A](shards) with Set[A]{
  /** The shards.  This ShardedSet object represents the union of £sets£. */ 
  protected val sets = 
    Array.fill(shards)(new scala.collection.mutable.HashSet[A])

  /** Does this contain £x£? */
  def contains(x: A): Boolean = {
    val s = shardFor(x); sets(s).synchronized{ sets(s).contains(x) }
  }

  /** Add £x£ to this.  Return £true£ if £x£ was not already in the set. */
  def add(x: A): Boolean = {
    val s = shardFor(x); sets(s).synchronized{ sets(s).add(x) }
  }

  /** Remove £x£ from this.  Return £true£ if £x£ was previously in the set. */
  def remove(x: A): Boolean = {
    val s = shardFor(x); sets(s).synchronized{ sets(s).remove(x) }
  }
}
\end{scala}
\caption{A sharded set implemented using JVM monitors.}
\label{fig:MonitorShardedSet}
\end{figure}

%%%%%

Figure~\ref{fig:MonitorShardedSet} gives a sharded set implemented using JVM
monitors.  Recall that the set is implemented in terms of some number of
smaller sets, known as shards.  Here we use |shards| |HashSet|s, stored in an
array |sets|.  A value~|x| is stored in the shard with index |shardFor(x)|
(defined in the class |Sharding|).  We want each shard to be independently
lockable, so operations on it are performed under mutual exclusion.  Thus we
need |shards| monitor objects.  We could
create |shards| new |AnyRef| objects, as we did in
Figure~\ref{fig:StripedCounterArray-JVM}.  However, a simpler method is to use
the |HashSet| objects themselves, as they are reference objects.  Hence each
operation on |set(s)| is protected by a |synchronized| block on |sets(s)|
itself.

However, it is not possible to re-implement the example from
Section~\ref{sec:linked-list-fine-grained} of a linked list using fine-grained
locking.  The extent of a |synchronized| block is defined by the brackets
around it.  Thus two different |synchronized| blocks by the same thread
satisfy the normal rules for bracketing: either they are disjoint, or one is
fully inside the other.  As noted in
Section~\ref{sec:linked-list-fine-grained}, the periods during which a thread
has two consecutive nodes locked does not satisfy this property.


